State-of-the-art LLMs are powered by scaling -- scaling model size, training tokens, and cluster size. It is economically infeasible to extensively tune hyperparameters for the largest runs. Instead, approximately optimal hyperparameters must be inferred or \textit{transferred} from smaller experiments. Hyperparameter transfer across model sizes has been studied in \citet{mup}. However, hyperparameter transfer across training tokens -- or token horizon -- has not been studied yet. To remedy this we conduct a large-scale empirical study on how optimal learning rate (LR) depends on the token horizon in LLM training. We first demonstrate that the optimal LR changes significantly with token horizon -- longer training necessitates smaller LR. Secondly, we demonstrate that the optimal LR follows a scaling law and that the optimal LR for longer horizons can be accurately estimated from shorter horizons via such scaling laws. We also provide a rule-of-thumb for transferring LR across token horizons with zero overhead over current practices. Lastly, we provide evidence that LLama-1 used too high LR, and thus argue that hyperparameter transfer across data size is an overlooked component of LLM training.
